%% kapa: exam=1 round=19 year=2008 period=1 form=A subject=civil_law subject_ko=민법 seq=20 number=20
%% source: Q-Net 감정평가사 기출문제 (cst00309) artlSeq=1001962 file=extracted/1008/1001962/2008년 제19회 감정평가사 1차시험 기출문제/1교시.hwp
%% conversion: kordoc 4.24.51 -> tools/gen1.py (auto); review=proofread-sonnet; flags=none
\begin{question}{20}
소멸시효의 중단에 관한 설명으로 틀린 것은? (다툼이 있는 경우에는 판례에 의함)
\begin{choices}
  \item 시효중단의 효력은 원칙적으로 당사자 및 그 승계인 사이에서만 효력이 있다.
  \item 재판상 청구가 있더라도 소의 기각이 있으면 시효중단의 효력이 있지만, 소의 각하가 있는 경우에는 시효중단의 효력이 없다.
  \item 근저당권설정등기청구의 소제기도 그 피담보채권이 될 채권에 대한 소멸시효 중단사유로 된다.
  \item 가압류의 집행행위가 시효의 이익을 받을 자에 대하여 행하여지지 않은 때에는 이를 그 자에게 통지한 후가 아니면 중단의 효력은 없다.
  \item 시효완성 후에 채무를 승인한 때에는 시효완성의 사실을 알고 그 이익을 포기한 것이라고 추정할 수 있다.
\end{choices}
\end{question}
\answer{2}
